\section{USAMO 2006}
        \begin{enumerate}
        	\item (\textit{USAMO 2006}) Let $p$ be a prime number and let $s$ be an integer with $0<s<p$. Prove that there exist integers $m$ and $n$ with $0<m<n<p$ and

 \[\left\lbrace\frac{sm}{p}\right\rbrace<\left\lbrace\frac{sn}{p}\right\rbrace<\frac{s}{p}\]
if and only if $s$ is not a divisor of $p-1$.


 Note: For $x$ a real number, let $\lfloor x\rfloor$ denote the greatest integer less than or equal to $x$, and let $\lbrace x\rbrace=x-\lfloor x\rfloor$ denote the fractional part of $x$.


 



\textbf{Solution 1}

We proceed by contradiction.  Assume that $s|(p-1)$.  Then for some positive integer $k$, $sk=p-1$.  The conditions given are equivalent to stating that $sm \bmod p < sn \bmod p< s\bmod p$.Now consider the following array modulo p:

 
 \[\begin{array}{ccccc} \mbox{Row 1}& s& 2s&\hdots& ks\\ \mbox{Row 2}& s-1& 2s-1& \hdots&ks-1\\ \vdots&\vdots&\vdots&\ddots&\vdots\\ \mbox{Row }s& 1& s+1& \hdots& (k-1)s+1 \end{array}\]

 Obviously, there are $s$ rows and $k$ columns. The first entry of each row is simply $((r-1)k+1)s\mod p$.  Since we wish for $sm\mod p$ and $sn\mod p$ to both be in the first column, while also satisfying the given conditions, we can easily see that $sm$ must be in a row $m_r$ with $m_r>n_r$, where $n_r$ denotes the row $sn \mod p$ is in.  However, since the values of each entry decreases while $((r-1)k+1)s$ keeps increasing, we can see that the condition can never be satisfied and thus, $s\not|(p-1).$

 To prove the other direction, let $sj+r=p-1$ for positive integers $j$ and $r$ with $j$ being the largest integer such that $sj<p-1$ and $r<s$. Since $s\not|(p-1)$, $1<s<p-1$.

 Note that for $0<l<j$, $ls\ge s\mod p$.  Thus, the first integer multiplied by s modulo p that will be less than s is $j+1.$ $(j+1)s\equiv s-r-1\mod p$.  Since $\lbrace s,2s,\hdots,js,(j+1)s,\hdots,xs,\hdots,(p-1)s \rbrace$ is a complete residue system mod p with the exception of the 0 term, we can find an $x>j+1$ such that $xs \equiv s-1 \mod p$.

 Thus, choose $m=j+1$ and $n=x$ to complete the proof.

 <i>Alternate solutions are always welcome. If you have a different, elegant solution to this problem, please \href{https://artofproblemsolving.com/wiki/index.php?title=2006\_USAMO\_Problems/Problem\_1&action=edit}{add it} to this page.</i>

 


	\item (\textit{USAMO 2006}) For a given positive integer $k$ find, in terms of $k$, the minimum value of $N$ for which there is a set of $2k+1$ distinct positive integers that has sum greater than $N$ but every subset of size $k$ has sum at most $N/2$.


 



\textbf{Solution 1}

Let one optimal set of integers be $\{a_1,\dots,a_{2k+1}\}$ with $a_1 > a_2 > \cdots > a_{2k+1} > 0$.

 The two conditions can now be rewritten as $a_1+\cdots + a_k \leq N/2$ and $a_1+\cdots +a_{2k+1} > N$.Subtracting, we get that $a_{k+1}+\cdots + a_{2k+1} > N/2$, and hence $a_{k+1}+\cdots + a_{2k+1} > a_1+\cdots + a_k$.In words, the sum of the $k+1$ smallest numbers must exceed the sum of the $k$ largest ones.

 Let $a_{k+1}=C$. As all the numbers are distinct integers, we must have $\forall i \in\{1,\dots,k\}:~ a_{k+1-i} \geq C+i$, and also $\forall i \in\{1,\dots,k\}:~ a_{k+1+i} \leq C-i$.

 Thus we get that $a_1+\cdots + a_k \geq kC + \dfrac{k(k+1)}2$, and $a_{k+1}+\cdots + a_{2k+1} \leq (k+1)C - \dfrac{k(k+1)}2$.

 As we want the second sum to be larger, clearly we must have $(k+1)C - \dfrac{k(k+1)}2 > kC + \dfrac{k(k+1)}2$.This simplifies to $C > k(k+1)$.

 Hence we get that:

 
 \begin{align*} N & \geq 2(a_1+\cdots + a_k) \\   & \geq 2\left( kC + \dfrac{k(k+1)}2 \right) \\   & = 2kC + k(k+1) \\   & \geq 2k(k^2+k+1) + k(k+1) \\   & = 2k^3 + 3k^2 + 3k \end{align*}


 On the other hand, for the set $\{ k^2+k+1+i ~|~ i\in\{-k,\dots,k\} \, \}$ the sum of the largest $k$ elements is exactly $k^3 + k^2 + k + \dfrac{k(k+1)}2$, and the sum of the entire set is $(k^2+k+1)(2k+1) = 2k^3 + 3k^2 + 3k + 1$, which is more than twice the sum of the largest set.

 Hence the smallest possible $N$ is $\boxed{ N = 2k^3 + 3k^2 + 3k }$.

 <i>Alternate solutions are always welcome. If you have a different, elegant solution to this problem, please \href{https://artofproblemsolving.com/wiki/index.php?title=2006\_USAMO\_Problems/Problem\_2&action=edit}{add it} to this page.</i>

 


	\item (\textit{USAMO 2006}) For integral $m$, let $p(m)$ be the greatest prime divisor of $m$. By convention, we set $p(\pm1)=1$ and $p(0)=\infty$. Find all polynomials $f$ with integer coefficients such that the sequence $\lbrace p(f(n^2))-2n\rbrace_{n\ge0}$ is bounded above. (In particular, this requires $f(n^2)\neq0$ for $n\ge0$.)


 



\textbf{Solution 1}

Let $f(x)$ be a non-constant polynomial in $x$ of degree $d$ withinteger coefficients, suppose further that no prime divides all thecoefficients of $f$ (otherwise consider the polynomial obtained bydividing $f$ by the gcd of its coefficients). We further normalize$f$ by multiplying by $-1$, if necessary, to ensure that theleading coefficient (of $x^d$) is positive.

 Let $g(n) = f(n^2)$, then $g(n)$ is a polynomial of degree $2$ ormore and $g(n) = g(-n)$. Let $g_1, \ldots, g_k$ be the factorizationof $g$ into irreducible factors with positive leading coefficients.Such a factorization is unique. Let $d(g_i)$ denote the degree of$g_i$. Since $g(-n) = g(n)$ the factors $g_i$ are either evenfunctions of $n$ or come in pairs $(g_i, h_i)$ with $g_i(-n) = (-1)^{d(g_i)} h_i(n)$.

 Let $P(0) = \infty$, $P(\pm 1) = 1$. For any other integer $m$ let$P(m)$ be the largest prime factor of $m$.

 Suppose that for some finite constant $C$ and all $n \ge 0$ we have$P(g(n)) - 2n < C$. Since the polynomials $g_i$ divide $g$, thesame must be true for each of the irreducible polynomials $g_i$.

 A theorem of T. Nagell implies that if $d(g_i) \ge 2$ the ratio$P(g_i(n))/n$ is unbounded for large values of $n$. Since in our case the $P(g_i(n))/n$ is asymptotically bounded aboveby $2$ for large $n$, we conclude that all the irreduciblefactors $g_i$ are linear. Since linear polynomials are not evenfunctions of $n$, they must occur in pairs $g_i(n) = a_in + b_i$,$h_i(n) = a_in - b_i$. Without loss of generality, $b_i \ge 0$.Since the coefficients of $f$ are relatively prime, so are $a_i$and $b_i$, and since $P(0) = \infty$, neither polynomial can haveany non-negative integer roots, so $a_i > 1$ and thus $b_i > 0$.

 On the other hand, by Dirichlet's theorem, $a_i \le 2$, sinceotherwise the sequence $a_in + b_i$ would yield infinitely manyprime values with $P(g_i(n)) = a_in + b_i \ge 3n.$ So $a_i = 2$ andtherefore $b_i$ is a positive odd integer. Setting $b_i = 2c_i + 1$, clearly $P(g_i(n)) - 2n < 2c_i + 2$.  Since this holds for eachfactor $g_i$, it is true for the product $g$ of all the factorswith the bound determined by the factor with the largest value of$c_i$.

 Therefore, for suitable non-negative integers $c_i$, $g(n)$ is aproduct of polynomials of the form $4n^2 - (2c_i + 1)^2$.  Now,since $g(n) = f(n^2)$, we conclude that $f(n)$ is a product oflinear factors of the form $4n - (2c_i + 1)^2$.

 Since we restricted ourselves to non-constant polynomials with relatively prime coefficients, we can now relax this condition andadmit a possibly empty list of linear factors as well as an arbitrarynon-zero integer multiple $M$. Thus for a suitable non-zero integer$M$ and $k \ge 0$ non-negative integers $c_i$, we have: 
 \[f(n) = M \cdot \prod_{i=1}^k (4n - (2c_i + 1)^2)\]

 



\textbf{Solution 2}

The polynomial $f$ has the required properties if and only if
 \[f(x) = c(4x - a_1^2)(4x - a_2^2)\cdots (4x - a_k^2),\qquad\qquad (*)\]where $a_1, a_2, \ldots, a_k$ are odd positive integers and $c$ is a nonzero integer. It is straightforward to verify that polynomials given by $(*)$ have the required property. If $p$ is a prime divisor of $f(n^2)$ but not of $c$, then $p|(2n - a_j)$ or $p|(2n + a_j)$ for some $j\leq k$. Hence $p - 2n\leq \max\{a_1, a_2, \ldots, a_k\}$. The prime divisors of $c$ form a finite set and do not affect whether or not the given sequence is bounded above. The rest of the proof is devoted to showing that any $f$ for which $\{p(f(n^2)) - 2n\}_{n\geq 0}$ is bounded above is given by $(*)$.

 Let $\mathbb{Z}[x]$ denote the set of all polynomials with integer coefficients. Given $f\in\mathbb{Z}[x]$, let $\mathcal{P}(f)$ denote the set of those primes that divide at least one of the numbers in the sequence $\{f(n)\}_{n\geq 0}$. The solution is based on the following lemma.

 \textbf{Lemma.} If $f\in\mathbb{Z}[x]$ is a nonconstant polynomial then $\mathcal{P}(f)$ is infinite.

 \textit{Proof.} Repeated use will be made of the following basic fact: if $a$ and $b$ are distinct integers and $f\in\mathbb{Z}[x]$, then $a - b$ divides $f(a) - f(b)$. If $f(0) = 0$, then $p$ divides $f(p)$ for every prime $p$, so $\mathcal{P}(f)$ is infinite. If $f(0) = 1$, then every prime divisor $p$ of $f(n!)$ satisfies $p > n$. Otherwise $p$ divides $n!$, which in turn divides $f(n!) - f(0) = f(n!) - 1$. This yields $p|1$, which is false. Hence $f(0) = 1$ implies that $\mathcal{P}(f)$ is infinite. To complete the proof, set $g(x) = f(f(0)x)/f(0)$ and observe that $g\in\mathcal{Z}[x]$ and $g(0) = 1$. The preceding argument shows that $\mathcal{P}(g)$ is infinite, and it follows that $\mathcal{P}(f)$ is infinite. $\blacksquare$

 Suppose $f\in\mathbb{Z}[x]$ is nonconstant and there exists a number $M$ such that $p(f(n^2)) - 2n\leq M$ for all $n\geq 0$. Application of the lemma to $f(x^2)$ shows that there is an infinite sequence of distinct primes $\{p_j\}$ and a corresponding infinite sequence of nonnegative integers $\{k_j\}$ such that $p_j|f(k_j)^2$ for all $j\geq 1$. Consider the sequence $\{r_j\}$ where $r_j = \min\{k_j\pmod{p_j}, p_j - k_j\pmod{p_j}\}$. Then $0\leq r_j\leq (p_j - 1)/2$ and $p_j|f(r_j)^2$. Hence $2r_j + 1\leq p_j\leq p(f(r_j^2))\leq M + 2r_j$, so $1\leq p_j - 2r_\leq M$ for all $j\geq 1$. It follows that there is an integer $a_1$ such that $1\leq a_1\leq M$ and $a_1 = p_j - 2r_j$ for infinitely many $j$. Let $m = \deg f$. Then $p_j|4^mf(((p_j - a_1)/2)^2)$ and $4^mf(((x - a_1)/2)^2)\in\mathbb{Z}[x]$. Consequently, $p_j|f((a_1/2)^2)$ for infinitely many $j$, which shows that $(a_1/2)^2$ is a zero of $f$. Since $f(n^2)\leq 0$ for $n\geq 0$, $a_1$ must be odd. Then $f(x) = (4x - a_1)^2g(x)$, where $g\in\mathbb{Z}[x]$. (See the note below.) Observe that $\{p(g(n^2)) - 2n\}_{n\geq 0}$ must be bounded above. If $g$ is constant, we are done. If $g$ is nonconstant, the argument can be repeated to show that $f$ is given by $(*)$.

 \textit{Note.} The step that gives $f(x) = (4x - a_1^2)g(x)$ where $g\in\mathbb{Z}[x]$ follows immediately using a lemma of Gauss. The use of such an advanced result can be avoided by first writing $f(x) = r(4x - a_1^2)g(x)$ where $r$ is rational and $g\in\mathbb{Z}[x]$. Then continuation gives $f(x) = c(4x - a_1^2)\cdots (4x - a_k^2)$ where $c$ is rational and the $a_i$ are odd. Consideration of the leading coefficient shows that the denominator of $c$ is $2^s$ for some $s\geq 0$ and consideration of the constant term shows that the denominator is odd. Hence $c$ is an integer.

 <i>Alternate solutions are always welcome. If you have a different, elegant solution to this problem, please \href{https://artofproblemsolving.com/wiki/index.php?title=2006\_USAMO\_Problems/Problem\_3&action=edit}{add it} to this page.</i>

 


	\item (\textit{USAMO 2006}) Find all positive integers $n$ for which there exists an integer $k\ge 2$ and positive rationals $a_1,a_2,\ldots a_k$ satisfying $a_1+a_2+\ldots+a_k=a_1\cdot a_2\cdots a_k=n$.


 



\textbf{Solution 1}

First, consider composite numbers.  We can then factor $n$ into $p_1p_2.$  It is easy to see that $p_1+p_2\le n$, and thus, we can add $(n-p_1-p_2)$ 1s in order to achieve a sum and product of $n$.  For $p_1+p_2=n$, which is only possible in one case, $n=4$, we consider $p_1=p_2=2$.  

 Secondly, let $n$ be a prime.  Then we can find the following procedure: Let $a_1=\frac{n}{2}, a_2=4, a_3=\frac{1}{2}$ and let the rest of the $a_k$ be 1.  The only numbers we now need to check are those such that $\frac{n}{2}+4+\frac{1}{2}>n\Longrightarrow n<9$.  Thus, we need to check for $n=1,2,3,5,7$.  One is included because it is neither prime nor composite.

 For $n=1$, consider $a_1a_2\hdots a_k=1$.  Then by AM-GM, $a_1+a_2+\hdots+a_k\ge k\sqrt[k]{1}>1$ for $k\ge 2$.  Thus, $n=1$ is impossible.

 If $n=2$, once again consider $a_1a_2\hdots a_k=2$.  Similar to the above, $a_1+a_2+\hdots\ge k\sqrt[k]{2}>2$ for $k\ge 2$ since $\sqrt[k]{2}>1$ and $k>2$. Obviously, $n=2$ is then impossible.

 If $n=3$, let $a_1a_2\hdots a_k=3$.  Again, $a_1+a_2+\hdots\ge k\sqrt[k]{3}>3$.  This is obvious for $k\ge 3$.  Now consider $k=2$.  Then $2\sqrt{3}\approx 3.4$ is obviously greater than $3$.  Thus, $n=3$ is impossible.

 If $n=5$, proceed as above and consider $k=2$.  Then $a_1+a_2=5$ and $a_1a_2=5$.  However, we then come to the quadratic $a_1^2-5a_1+5=0 \Longrightarrow a_1=\frac{5\pm\sqrt{5}}{2}$, which is not rational.  For $k=3$ and $k=4$ we note that $\sqrt[3]{5}>\frac{5}{3}$ and $\sqrt[4]{5}>\frac{5}{4}$.  This is trivial to prove.  If $k\ge 5$, it is obviously impossible, and thus $n=5$ does not work.

 The last case, where $n=7$, is possible using the following three numbers.  $a_1=\frac{9}{2}, a_2=\frac{4}{3}, a_3=\frac{7}{6}$ shows that $n=7$ is possible.

 Hence, $n$ can be any positive integer greater than $3$ with the exclusion of $5$.

 <i>Alternate solutions are always welcome. If you have a different, elegant solution to this problem, please \href{https://artofproblemsolving.com/wiki/index.php?title=2006\_USAMO\_Problems/Problem\_4&action=edit}{add it} to this page.</i>

 


	\item (\textit{USAMO 2006}) A mathematical frog jumps along the number line. The frog starts at 1, and jumps according to the following rule: if the frog is at integer $n$, then it can jump either to $n+1$ or to $n+2^{m_n+1}$ where $2^{m_n}$ is the largest power of 2 that is a factor of $n$. Show that if $k\ge2$ is a positive integer and $i$ is a nonnegative integer, then the minimum number of jumps needed to reach $2^ik$ is greater than the minimum number of jumps needed to reach $2^i$.


 



\textbf{Solution 1}

For $i\geq 0$ and $k\geq 1$, let $x_{i,k}$ denote the minimum number of jumps needed to reach the integer $n_{i,k} = 2^i k$. We must prove that
 \[x_{i,k} > x_{i,1}\qquad\qquad (1)\]for all $i\geq 0$ and $k\geq 2$. We prove this using the method of descent.

 First note that $(1)$ holds for $i = 0$ and all $k\geq 2$, because it takes 0 jumps to reach the starting value $n_{0,1} = 1$, and at least one jump to reach $n_{0,k} = k\geq 2$. Now assume that $(1)$ is not true for all choices of $i$ and $k$. Let $i_0$ be the minimal value of $i$ for which $(1)$ fails for some $k$, let $k_0$ be the minimal value of $k > 1$ for which $x_{i_0,k}\leq x_{i_0,1}$. Then it must be the case that $i_0\geq 1$ and $k_0\geq 2$.

 Let $J_{i_0,k_0}$ be a shortest sequence of $x_{i_0,k_0} + 1$ integers that the frog occupies in jumping from 1 to $2^{i_0} k_0$. The length of each jump, that is, the difference between consecutive integers in $J_{i_0,k_0}$, is either 1 or a positive integer power of 2. The sequence $J_{i_0,k_0}$ cannot contain $2^{t_0}$ because it takes more jumps to reach $2^{t_0} k_0$ than it does to reach $2^{t_0}$. Let $2^{M+1}$, $M\geq 0$ be the length of the longest jump made in generating $J_{i_0,k_0}$. Such a jump can only be made from a number that is divisible by $2^M$ (and by no higher power of 2). Thus we must have $M < i_0$, since otherwise a number divisible by $2^{i_0}$ is visited before $2^{i_0} k_0$ is reached, contradicting the definition of $k_0$.

 Let $2^{m+1}$ be the length of the jump when the frog jumps over $2^{i_0}$. If this jump starts at $2^m(2t - 1)$ for some positive integer $t$, then it will end at $2^m(2t - 1) + 2^{m+1} = 2^m(2t + 1)$. Since it goes over $2^{i_0}$ we see $2^m(2t - 1) < 2^{i_0} < 2^m(2t + 1)$ or $(2^{i_0-m} - 1)/2 < t < (2^{i_0-m} + 1)/2$. Thus $t = 2^{i_0-m-1}$ and the jump over $2^{i_0}$ is from $2^m(2^{i_0-m} - 1) = 2^{i_0} - 2^m$ to $2^m(2^{i_0-m}+1) = 2^{i_0} + 2^m$.

 Considering the jumps that generate $J_{i_0,k_0}$, let $N_1$ be the number of jumps from 1 to $2^{i_0} + 2^m$, and let $N_2$ be the number of jumps from $2^{i_0} + 2^m$ to $2^{i_0}k$. By definition of $i_0$, it follows that $2^m$ can be reached from 1 in less than $N_1$ jumps. On the other hand, because $m < i_0$, the number $2^{i_0}(k_0 - 1)$ can be reached from $2^m$ in exactly $N_2$ jumps by using the same jump length sequence as in jumping from $2^m + 2^{i_0}$ to $2^{i_0} k_0 = 2^{i_0}(k_0 - 1) + 2^{i_0}$. The key point here is that the shift by $2^{i_0}$ does not affect any of divisibility conditions needed to make jumps of the same length. In particular, with the exception of the last entry, $2^{i_0} k_0$, all of the elements of $J_{i_0,k_0}$ are of the form $2^p(2t + 1)$ with $p < i_0$, again because of the definition of $k_0$. Because $2^p(2t + 1) - 2^{i_0} = 2^p(2t - 2^{i_0-p} + 1)$ and the number $2t + 2^{i_0-p} + 1$ is odd, a jump of size $2^{p+1}$ can be made from $2^p(2t + 1) - 2^{i_0}$ just as it can be made from $2^p(2t + 1)$.

 Thus the frog can reach $2^m$ from 1 in less than $N_1$ jumps, and can then reach $2^{i_0}(k_0 - 1)$ from $2^m$ in $N_2$ jumps. Hence the frog can reach $2^{i_0}(k_0 - 1)$ from 1 in less than $N_1 + N_2$ jumps, that is, in fewer jumps than needed to get to $2^{i_0} k_0$ and hence in fewer jumps than required to get to $2^{i_0}$. This contradicts the definition of $k_0$.

 



\textbf{Solution 2}

Suppose $x_0 = 1, x_1, \ldots, x_t = 2^i k$ are the integers visited by the frog on his trip from 1 to $2^i k$, $k\geq 2$. Let $s_j = x_j - x_{j-1}$ be the jump sizes. Define a reduced path $y_j$ inductively by
 \[y_j = \left\{\begin{array}{ll} y_{j-1} + s_j & \text{if }y_{j-1} + s_j\leq 2^i, \\ y_{j-1} & \text{otherwise.} \end{array}\right.\]Say a jump $s_j$ is deleted in the second case. We will show that the distinct integers among the $y_j$ give a shorter path from 1 to $2^i$. Clearly $y_j\leq 2^i$ for all $j$. Suppose $2^i - 2^{r+1} < y_j\leq 2^i - 2^r$ for some $0\leq r\leq i - 1$. Then every deleted jump before $y_j$ must have length greater than $2^r$, hence must be a multiple of $2^{r+1}$. Thus $y_j\equiv x_j\pmod{2^{r+1}}$. If $y_{j+1} > y_j$, then either $s_{j+1} = 1$ (in which case this is a valid jump) or $s_{j+1}/2 = 2^m$ is the exact power of 2 dividing $x_j$. In the second case, since $2^r\geq s_{j+1} > 2^m$, the congruence says $2^m$ is also the exact power of 2 dividing $y_j$, thus again this is a valid jump. Thus the distinct $y_j$ form a valid path for the frog. If $j = t$ the congruence gives $y_t\equiv x_t\equiv 0\pmod{2^{r+1}}$, but this is impossible for $2^i - 2^{r+1} < y_t\leq 2^i - 2^r$. Hence we see $y_t = 2^i$, that is, the reduced path ends at $2^i$. Finally since the reduced path ends at $2^i < 2^i k$ at least one jump must have been deleted and it is strictly shorter than the original path.

 <i>Alternate solutions are always welcome. If you have a different, elegant solution to this problem, please \href{https://artofproblemsolving.com/wiki/index.php?title=2006\_USAMO\_Problems/Problem\_5&action=edit}{add it} to this page.</i>

 


	\item (\textit{USAMO 2006}) Let $ABCD$ be a quadrilateral, and let $E$ and $F$ be points on sides $AD$ and $BC$, respectively, such that $AE/ED=BF/FC$.  Ray $FE$ meets rays $BA$ and $CD$ at $S$ and $T$ respectively. Prove that the circumcircles of triangles $SAE$, $SBF$, $TCF$, and $TDE$ pass through a common point.


 



\textbf{Solution 1}

Let the intersection of the circumcircles of $SAE$ and $SBF$ be $X$, and let the intersection of the circumcircles of $TCF$ and $TDE$ be $Y$.

 $BXF=BSF=AXE$ because $BSF$ tends both arcs $AE$ and $BF$.$BFX=XSB=XEA$ because $XSB$ tends both arcs $XA$ and $XB$.Thus, $XAE\sim XBF$ by AA similarity, and $X$ is the center of spiral similarity for $A,E,B,$ and $F$.$FYC=FTC=EYD$ because $FTC$ tends both arcs $ED$ and $FC$.$FCY=FTY=EDY$ because $FTY$ tends both arcs $YF$ and $YE$.Thus, $YED\sim YFC$ by AA similarity, and $Y$ is the center of spiral similarity for $E,D,F,$ and $C$.

 From the similarity, we have that $XE/XF=AE/BF$. But we are given $ED/AE=CF/BF$, so multiplying the 2 equations together gets us $ED/FC=XE/XF$. $DEX,CFX$ are the supplements of $AEX, BFX$, which are congruent, so $DEX=CFX$, and so $XED\sim XFC$ by SAS similarity, and so $X$ is also the center of spiral similarity for $E,D,F,$ and $C$. Thus, $X$ and $Y$ are the same point, which all the circumcircles pass through, and so the statement is true.

 



\textbf{Solution 2}

We will give a solution using complex coordinates. The first step is the following lemma.

 \textbf{Lemma.} Suppose $s$ and $t$ are real numbers and $x$, $y$ and $z$ are complex. The circle in the complex plane passing through $x$, $x + ty$ and $x + (s + t)z$ also passes through the point $x + syz/(y - z)$, independent of $t$.

 \textit{Proof.} Four points $z_1$, $z_2$, $z_3$ and $z_4$ in the complex plane lie on a circle if and only if the cross-ratio
 \[cr(z_1, z_2, z_3, z_4) = \frac{(z_1 - z_3)(z_2 - z_4)}{(z_1 - z_4)(z_2 - z_3)}\]is real. Since we compute
 \[cr(x, x + ty, x + (s + t)z, x + syz/(y - z)) = \frac{s + t}{s}\]the given points are on a circle. $\blacksquare$

 Lay down complex coordinates with $S = 0$ and $E$ and $F$ on the positive real axis. Then there are real $r_1$, $r_2$ and $R$ with $B = r_1A$, $F = r_2E$ and $D = E + R(A - E)$ and hence $AE/ED = BF/FC$ gives
 \[C = F + R(B - F) = r_2(1 - R)E + r_1RA.\]The line $CD$ consists of all points of the form $sC + (1 - s)D$ for real $s$. Since $T$ lies on this line and has zero imaginary part, we see from $\text{Im}(sC + (1 - s)D) = (sr_1R + (1 - s)R)\text{Im}(A)$ that it corresponds to $s = -1/(r_1 - 1)$. Thus
 \[T = \frac{r_1D - C}{r_1 - 1} = \frac{(r_2 - r_1)(R - 1)E}{r_1 - 1}.\]Apply the lemma with $x = E$, $y = A - E$, $z = (r_2 - r_1)E/(r_1 - 1)$, and $s = (r_2 - 1)(r_1 - r_2)$. Setting $t = 1$ gives
 \[(x, x + y, x + (s + 1)z) = (E, A, S = 0)\]and setting $t = R$ gives
 \[(x, x + Ry, x + (s + R)z) = (E, D, T).\]Therefore the circumcircles to $SAE$ and $TDE$ meet at
 \[x + \frac{syz}{y - z} = \frac{AE(r_1 - r_2)}{(1 - r_1)E - (1 - r_2)A} = \frac{AF - BE}{A + F - B - E}.\]This last expression is invariant under simultaneously interchanging $A$ and $B$ and interchanging $E$ and $F$. Therefore it is also the intersection of the circumcircles of $SBF$ and $TCF$.

 



\textbf{Solution 3}

Let $M$ be the Miquel point of $ABCD$; then $M$ is the center of the spiral similarity that takes $AD$ to $BC$. Because $\frac{AE}{ED} = \frac{BF}{FC}$, the same spiral similarity also takes $E$ to $F$, so $M$ is the center of the spiral similarity that maps $AE$ to $BF$ and $ED$ to $FC$. Then it is obvious that the circumcircles of $SAE$, $SBF$, $TCF$, and $TDE$ pass through $M$.

 <i>Alternate solutions are always welcome. If you have a different, elegant solution to this problem, please \href{https://artofproblemsolving.com/wiki/index.php?title=2006\_USAMO\_Problems/Problem\_6&action=edit}{add it} to this page.</i>

 


\end{enumerate}